\subsection{Platform Selection}

Currently, there are three main application platforms: web applications, native mobile applications, and hybrid mobile applications. Web applications are accessible from any device through a browser by entering the URL. Native mobile applications are installed on a mobile phone and are specific to a single operating system, either Android or iOS. Hybrid mobile applications combine the characteristics of both web and native platforms.

A native application requires deep specialization for both operating systems, demanding higher development costs. Its adoption may also be slower, as users must install the application before using it. While mobile phones are convenient for uploading the expenses via the camera, performing the subsequent AI-extraction review and report creation is much more comfortable on a larger screen.

Conversely, a web application offers more practical advantages. A web application can be accessed from any device without installation. Development is much faster, a single codebase implementation is more than enough, because it is valid for all users. Ongoing maintenance and updates are instant after simply publishing the new version on the server. 


Given the primary use case of creating expense reports at a desk and the need for universal accessibility without installation barriers, the \textbf{web platform} is selected. In terms of cost and development effort, this one is also the most suitable option.
